\FloatBarrier

\section{Scope and geometry changes}

The question is how much the optimized LTCC CAD implementation changes loading cost, memory use,
particle-transport cost, and geometry diagnostics. This is a software and geometry study, not a background-rate
prediction or a measurement of LTCC photon yield. No experimental data were collected.

The original data are taken from main revision ea3625f. The port uses the LTCC CAD changes in meson revision
4c7c6af. The later PMT absorption change in revision 3472ede is excluded, so PMT optical properties remain the
same in the compared geometries. The unrelated Meson build-system changes are also excluded; the same installed
executable runs every case. The source revisions are recorded in Refs.~\cite{baseline,optimized}.

The primary comparison restores all ten files from main's cadTempRemoved directory: S1-BW, S1-LW, S1-RW,
S1-TRB, NFrame, FrontPlate, Mount, Nose, BottomPlate, and Epoxy. Source volumes are placed in sector 3 and the
disabled frame and six-part nose copies are enabled in the other five sectors. The benchmark preserves the
legacy placement formulas and material assignments from ltcc\_frame.pl: source frame volumes are aluminum,
copied frames are stainless steel, nose pieces are aluminum, and Epoxy is G4\_CR39. No STL facet is repaired or
removed in the original baseline. This gives 90 frame/nose physical placements, compared with 24 in the
consolidated implementation, excluding the Winston cones.

The previous production configuration had this hardware disabled. Comparing it directly with the complete meson
model changes the amount of simulated hardware and answers a different question. Those initial measurements are
retained as supplementary context; the headline optimization results in this note compare the complete original
and consolidated models.

The consolidated implementation uses four hardware source volumes: backwall, rightwall, leftwall, and nose.
Sector 3 contains the source volumes; enabled sectors reuse them through CopyOf placements inside their LTCC
mothers. The port restores walls and nose hardware that were disabled in the original geometry, changes the
frame hierarchy and some materials, and replaces the previous chain of mother-volume subtractions with a
fitted, planar G4Trap intersection. The unmodified engineering drawings remain archived in the geometry source
repository. The restored original baseline keeps the legacy mothers and hierarchy unchanged.

The consolidated nose uses the implementation's effective material at density 3.86 g/cm\textsuperscript{3},
with mass fractions 0.9464 stainless steel, 0.0408 carbon, 0.0043 hydrogen, and 0.0085 oxygen. This study does
not establish material-budget or detector-response equivalence to the original model.

Three cases separate part of the optimization: Original uses the old data; Cones only replaces only WC\_S,
WC\_M, and WC\_L while retaining the restored original frame, mothers, and materials; Full port uses the new
geometry and all its optimized meshes. Cones only is a diagnostic configuration, not an implementation proposed
for deployment. Frame consolidation and mother changes remain combined in the Full port case and cannot be
assigned independent transport speedups from this comparison.

Figures~\ref{fig:stl-hardware}--\ref{fig:stl-cones} compare renderings of the original and optimized STL files
directly. The original includes all restored CAD hardware. Paired panels use identical orthographic cameras,
common bounds, and the same scale; no mesh repair, smoothing, or additional decimation is applied. The hardware
view applies the sector-3 placements, including the sector mother transform for the optimized model. The
resulting hardware transforms are checked against the restored original placements. Nose and cone close-ups use
their stored STL coordinate systems. Colors identify functional hardware groups, not materials. The images
expose simplification and changes in shape; they do not establish volume, material-budget, or detector-response
equivalence.

\begin{figure}[htbp]
\centering
\includegraphics[width=\textwidth]{figures/hardware_stl_comparison.pdf}
\caption{Sector-3 STL frame and nose hardware: complete original on the left and optimized on the right. Upper
panels show shaded front oblique views; lower panels show rear oblique views with triangle edges. The same
camera and scale are used within each row. The 15 original meshes contain 124,652 triangles; the four optimized
meshes contain 3,646. Winston cones, native LTCC mirrors, PMTs, mothers, and other sector copies are omitted.}
\label{fig:stl-hardware}
\end{figure}

\begin{figure}[p]
\centering
\includegraphics[width=\textwidth]{figures/nose_stl_comparison.pdf}
\caption{Nose consolidation in the stored engineering coordinates, with identical camera and scale: six
original meshes and 12,012 triangles on the left; one optimized mesh and 1,986 triangles on the right. Triangle
edges are shown; simplified details and differences in the rendered envelope are visible.}
\label{fig:stl-nose}
\end{figure}

\begin{figure}[p]
\centering
\includegraphics[width=\textwidth]{figures/cones_stl_comparison.pdf}
\caption{Winston-cone meshes with triangle edges: WC\_S, WC\_M, and WC\_L from top to bottom; original on the
left and optimized on the right. The camera, bounds, and scale match within each row. Original triangle counts
are 4,312, 3,076, and 2,578; each optimized mesh contains 1,500 triangles. These are individual source meshes,
not all physical cone placements.}
\label{fig:stl-cones}
\end{figure}
